\chapter{Prototype Development}

To be able to prove the proposed solution is viable and will successfully address the gap identified, a prototype needs to be developed and afterwards validated. % TODO add a bit more here

% \section{Technology Selection}
% Firstly some decisions need to be made to choose the most appropriate technologies to support this work.

\section{Runtime Environment}
Out of all \acp{RTE} compatible with \textit{IEC 61499}, this prototype will focus on supporting \textit{DINASORE} (Dynamic INtelligent Architecture for Software and MOdular REconfiguration). It is much less complex than others, making it ideal for a proof of concept \ac{IDE}. As an open-source project, \textit{DINASORE} can be freely inspected and studied, which is particularly valuable in a research and prototyping context where flexibility is prioritised over real-time performance guarantees. This ease of study rules out the use of commercial \acp{RTE} like \textit{EcoStruxure Automation Expert's "dPAC"} and \textit{ISaGRAF Runtime}. Other open-source \acp{RTE} were considered, but none was as well-suited as \textit{DINASORE}. \textit{4diac FORTE} for example was way too extensive and complicated to be able to easily study it's source code.
% TODO maybe talk about the use of python both in dinasores code and the fbs code

\begin{figure}[h!]
  
\includegraphics[width=0.5\linewidth]{dinasore-logo.png}
  \caption{\textit{DINASORE} logo}
  \label{fig:dinasore-logo}
\end{figure}


\section{Server}
The design begins by setting up a server that will serve the \ac{IDE} interface as a webpage, store and manage all the data and handle communication with the \acp{RTE}. Taking into account the features we plan on implementing, \textit{JavaScript} running on \textit{Node.js} were chosen due to the extensive library ecosystem. For real-time collaboration we will need a library to run in both the server and the frontend, since \textit{JavaScript} runs natively in the browser, it further reinforces that it should also be used in the server. The library choice will be explained in depth in \Cref{sec:real-time-collaboration}, proving the previous train of thought.

Being the middleware, there are 4 different types of network communications that the server will need to handle, each with it's own requirements:

\begin{itemize}
  \item \textbf{Serve webpage}: Communicates via \textit{GET} requests to deliver the \textit{html}, \textit{css} and \textit{js} files that make up the \ac{IDE} interface.
  \item \textbf{Real-Time Collaboration}: Constant atomic updates via a library.
  \item \textbf{Requested deployment}: A team member sends a \textit{POST} request, that tells the server to deploy.
  \item \textbf{Deployment sockets}: Communicates via sockets with each \ac{RTE} to send commands and receive the corresponding responses.
\end{itemize}

\begin{figure}[h!]
  \includegraphics[width=1\linewidth]{ide-server-rte.jpg}
  \caption{Communication through the Middleware}
  \label{fig:ide-server-rte}
\end{figure}

\section{Frontend}
The frontend is the \ac{UI} that will be used to create and modify the \ac{FBD} programs. It consists of a 3 views.

\subsection{Devices Configuration}

\subsection{Function Blocks Editor View}

\subsection{Graph View}



\section{Real-Time Collaboration} \label{sec:real-time-collaboration}


\section{Deployment}
\subsection{Deployment Console}
